% New English translation from checked cote 25, site 91 (PDF page 92).
% AI-translated by GPT-6 Astra (OpenAI, ChatGPT Pro mode).
\scanpage{25/91}
\mspage{85}

Of course, to justify in the \var{\xout{present}}{}{typewriter strike-out}{ega25_25_091-v0001} case of the present no. the passage
from generic \hand{$\eta$} to neighboring points of $\mathrm{Grass}^{m+1}(P)$, and also to be able,
where appropriate, to use for our purposes the general explanations of
7.1., we should consider the diagrams
\diagram{\[\begin{array}{ccc}
X & \hand{\longleftarrow} & \overset{\hand{\sim}}{X}(C) \\
\hand{\downarrow} & & \hand{\downarrow} \\
S & \hand{\longleftarrow} & \hand{Q}(C)
\end{array}\]}
obtained using the various $C \hand{\in} \mathrm{Grass}^{m+1}(S)$, and \var{\inserted[typed]{more generally}}{more generally}{insertion}{ega25_25_091-v0002} those obtained likewise
after base change $T \hand{\to} S$ for points \hand{$\unsure{\xi|\zeta}$} \hand{$\in$} $\mathrm{Grass}^{m+1}($\var{\over{P}{\unsure{T}}}{\unsure{T}}{writing over}{ega25_25_091-v0003}$)$:
\diagram{\[\begin{array}{ccc}
X_T & \hand{\longleftarrow} & \overset{\hand{\sim}}{X}(C_{\hand{\unsure{\xi|\zeta}}}) = \overset{\hand{\sim}}{X}_T(C_{\hand{\unsure{\xi|\zeta}}}) \\
\hand{\downarrow} & & \hand{\downarrow} \\
T & \hand{\longleftarrow} & \hand{Q}(C_{\hand{\unsure{\xi|\zeta}}})
\end{array}\]}
as obtained, by base change\hand{s} \hand{$\unsure{\xi|\zeta}$}: $T \hand{\longrightarrow} \mathrm{Grass}^{m+1}$\var{\deleted{$(P)$}}{}{deletion}{ega25_25_091-v0004} $= \underline{T}$,
from the \uline{universal} diagram (relative to the canonical point of $\mathrm{Grass}^{m+1}$ in $\underline{T}$):
\diagram{\[\begin{array}{ccc}
X_{\underline{T}} & \hand{\longleftarrow} & \overset{\hand{\sim}}{X}(\underline{C}) \\
\hand{\downarrow} & & \hand{\downarrow} \\
\underline{T} & \hand{\longleftarrow} & \hand{Q}(\underline{C})
\end{array}\]}
where $\underline{C}$ is the canonical linear subvariety of $P_{\underline{T}}$. Then above
$\overset{\hand{\sim}}{X}(C_{\hand{\unsure{\eta|y|\zeta}}}) \hand{\longrightarrow} \hand{Q}(C_{\hand{\unsure{\eta|y|\zeta}}})$ is precisely the morphism of generic fibers for
the $\underline{T}$-morphism $\overset{\hand{\sim}}{X}(\underline{C}) \hand{\longrightarrow} \hand{Q}(\underline{C})$ of this last diagram, and every constructible
property for \var{\xout{\illegible{1}}}{}{typewriter strike-out}{ega25_25_091-v0005} the morphism of generic fibers implies
the property for neighboring fibers. -- From the notational viewpoint,
$\hand{Q}$ should be regarded (and even introduced) as the name of the natural morphism \var{\inserted[typed]{of functors}}{of functors}{insertion}{ega25_25_091-v0006}
$\mathrm{Grass}^{m+1}(P) \hand{\longrightarrow}$ \var{\xout{Hilb}}{}{typewriter strike-out}{ega25_25_091-v0007} Subpreschemes $(\mathrm{Grass}^{m}(P))$.

15. \uline{Axiomatization of some of the preceding results}.

I am thinking above all of the results of nos. 2 to 8, which for the most part
\tnote{25-091-1}{The arguments S and P in the successive Grass expressions, and the subsequent deletion of (P), follow the source. Subpreschemes translates the literal functor name Sous-préschémas.}

\mspage{91}
